\documentclass{article}
\usepackage{fleqn}
\usepackage{epsf}
\usepackage[dvips]{color}
\usepackage{aima2e-slides}

\begin{document}

\begin{huge}
\titleslide{First-order logic}{Chapter 8}

\sf

%%%%%%%%%%%% Slide %%%%%%%%%%%%%%%%%%%%%%%%%%%%%%%%%%%%%%%%%%%%%%%%%%%
\heading{Phác thảo}

\blob Tại sao FOL?

\blob Cú pháp và ngữ nghĩa của FOL

\blob Vui với những câu nói

\blob Thế giới Wumpus trong FOL


%%%%%%%%%%%% Slide %%%%%%%%%%%%%%%%%%%%%%%%%%%%%%%%%%%%%%%%%%%%%%%%%%%
\heading{Ưu điểm và nhược điểm của logic mệnh đề}

\smiley\ Logic mệnh đề là \emph{declarative}: các đoạn cú pháp tương ứng với các sự kiện

\smiley\ Logic mệnh đề cho phép thông tin một phần/rời rạc/phủ định \al
  (không giống như hầu hết các cấu trúc dữ liệu và cơ sở dữ liệu)

\smiley\ Logic mệnh đề là \emph{thành phần}: \al
  ý nghĩa của \mat{$B_{1,1}\land P_{1,2}$} được bắt nguồn từ ý nghĩa của \mat{$B_{1,1}$} và của \mat{$P_{1,2}$}

\smiley\ Ý nghĩa trong logic mệnh đề là \emph{không phụ thuộc vào ngữ cảnh}\al
  (không giống như ngôn ngữ tự nhiên, ý nghĩa phụ thuộc vào ngữ cảnh)

\frowny\ Logic mệnh đề có khả năng biểu đạt rất hạn chế\al
  (không giống ngôn ngữ tự nhiên)\al
  Ví dụ: không thể nói `` hố gây gió ở các ô vuông liền kề ''\nl
    ngoại trừ việc viết một câu cho mỗi ô vuông

%%%%%%%%%%%% Slide %%%%%%%%%%%%%%%%%%%%%%%%%%%%%%%%%%%%%%%%%%%%%%%%%%%
\heading{Logic bậc nhất}

Trong khi logic mệnh đề giả định thế giới chứa \emph{sự kiện},\\
logic bậc nhất (như ngôn ngữ tự nhiên) giả định thế giới chứa đựng
\begin{itemize} 
\item \defn{Đối tượng}: con người, nhà cửa, con số, lý thuyết,
Ronald McDonald, màu sắc, trò chơi bóng chày, chiến tranh, thế kỷ $\ldots$ 
\item \defn{Quan hệ}: đỏ, tròn, giả, nguyên tố, nhiều tầng $\ldots$, \\
anh trai của, lớn hơn, bên trong, một phần của, có màu sắc, xảy ra sau, sở hữu, ở giữa, $\ldots$
\item \defn{Chức năng}: cha của, bạn thân, hiệp thứ ba, hiệp nữa
hơn, cuối $\ldots$
\end{itemize}

%%%%%%%%%%%% Slide %%%%%%%%%%%%%%%%%%%%%%%%%%%%%%%%%%%%%%%%%%%%%%%%%%%
\heading{Logic nói chung}

\begin{mytabular}{@{\extracolsep\fill}|@{\squad}l@{\quad}|l@{\quad}|l@{\quad}|}
\hline
\tabhead {\tf Ngôn ngữ (Language)} & {\tf Cam kết Bản thể học (Ontological Commitment)} & {\tf Cam kết Nhận thức học (Epistemological Commitment)} \\
\hline
\tabtop Logic mệnh đề (Propositional logic) & sự kiện & đúng/sai/không biết \\
Logic bậc một (First-order logic) & sự kiện, đối tượng, quan hệ & đúng/sai/không biết\\
Logic thời gian (Temporal logic) & sự kiện, đối tượng, quan hệ, thời gian & đúng/sai/không biết\\
Lý thuyết xác suất (Probability theory) & sự kiện & mức độ tin tưởng 0\ldots 1\\
\tabbot Logic mờ (Fuzzy logic) & mức độ đúng đắn & mức độ tin tưởng 0\ldots 1\\
\hline
\end{mytabular}





%%%%%%%%%%%% Slide %%%%%%%%%%%%%%%%%%%%%%%%%%%%%%%%%%%%%%%%%%%%%%%%%%%
\heading{Cú pháp của FOL: Các phần tử cơ bản}

\begin{tabular}{ll}
Hằng số & \mat{$KingJohn,\ 2,\ UCB,\ldots$} \\
Vị ngữ & \mat{$Brother,\ >,\ldots$} \\
Chức năng & \mat{$Sqrt,\ LeftLegOf,\ldots$} \\
Biến & \mat{$x,\ y,\ a,\ b,\ldots$} \\
Từ nối & \mat{$\land\ \lor\ \lnot\ \implies\ \lequiv$} \\
Bình đẳng & \mat{$=$} \\
Bộ định lượng & \mat{$\forall\ \exists$} \\
\end{tabular}



%%%%%%%%%%%% Slide %%%%%%%%%%%%%%%%%%%%%%%%%%%%%%%%%%%%%%%%%%%%%%%%%%%
\heading{Câu nguyên tử}

\begin{tabular}{rcl}
Câu nguyên tử & = & \mat{$predicate(term_1,\ldots,term_n)$} \\
                & & hoặc \mat{$term_1 = term_2$} \\
                & & \\
Thuật ngữ & = & \mat{$function(term_1,\ldots,term_n)$} \\
                & & hoặc \mat{$constant$} hoặc \mat{$variable$} \\
\end{tabular}

\begin{tabular}{ll}
Ví dụ: & \mat{$Brother(KingJohn,RichardTheLionheart)$}\\
      & \mat{$>(Length(LeftLegOf(Richard)),Length(LeftLegOf(KingJohn)))$}\\
\end{tabular}


%%%%%%%%%%%% Slide %%%%%%%%%%%%%%%%%%%%%%%%%%%%%%%%%%%%%%%%%%%%%%%%%%%
\heading{Câu phức tạp}

Câu phức được tạo thành từ câu nguyên tử bằng cách sử dụng từ nối
\mat{\[
\lnot S,\quad S_1\land S_2,\quad S_1 \lor S_2,\quad 
S_1 \implies S_2,\quad S_1 \lequiv S_2
\]}

\begin{tabular}{ll}
Ví dụ: & \mat{$Sibling(KingJohn,Richard) \implies Sibling(Richard,KingJohn)$} \\
     & \mat{${>}(1,2) \lor {\leq}(1,2)$} \\
     & \mat{${>}(1,2) \land \lnot {>}(1,2)$} \\
\end{tabular}



%%%%%%%%%%%% Slide %%%%%%%%%%%%%%%%%%%%%%%%%%%%%%%%%%%%%%%%%%%%%%%%%%%
\heading{Sự thật trong logic bậc nhất}

Các câu đều đúng đối với mô hình \defn{} và cách diễn giải \defn{}

Mô hình chứa các đối tượng $\geq 1$ (\defn{các phần tử miền }) và mối quan hệ giữa chúng

Giải thích chỉ định tham chiếu cho \al
\mat{constant symbols} $\rightarrow$ \note{đối tượng}\al
\mat{predicate symbols} $\rightarrow$ \note{quan hệ}\al
\mat{function symbols} $\rightarrow$ \note{quan hệ chức năng}

Một câu nguyên tử \mat{$predicate(term_1,\ldots,term_n)$} là đúng\\
nếu \note{đối tượng} được đề cập bởi \mat{$term_1,\ldots,term_n$}\\
nằm trong mối quan hệ \note{} được đề cập bởi \mat{$predicate$}



%%%%%%%%%%%% Slide %%%%%%%%%%%%%%%%%%%%%%%%%%%%%%%%%%%%%%%%%%%%%%%%%%%
\heading{Mô hình cho FOL: Ví dụ}

\vspace*{0.3in}

\epsfxsize=0,85\textwidth
\fig{\file{figures}{fol-model.ps}}



%%%%%%%%%%%% Slide %%%%%%%%%%%%%%%%%%%%%%%%%%%%%%%%%%%%%%%%%%%%%%%%%%%
\heading{Ví dụ về sự thật}

Hãy xem xét cách giải thích trong đó\\
\mat{$Richard$} $\rightarrow$ \note{Richard Tim Sư Tử}\\
\mat{$John$} $\rightarrow$ \note{Vua độc ác John}\\
\mat{$Brother$} $\rightarrow$ \note{Mối quan hệ anh em}

Theo cách giải thích này, \mat{$Brother(Richard,John)$} là đúng\\
chỉ trong trường hợp \note{Richard the Lionheart} và \note{Vua độc ác John}\\
nằm trong \note{mối quan hệ anh em} trong mô hình


%%%%%%%%%%%% Slide %%%%%%%%%%%%%%%%%%%%%%%%%%%%%%%%%%%%%%%%%%%%%%%%%%%
\heading{Mô hình cho FOL: Rất nhiều!}

Sự đòi hỏi trong logic mệnh đề có thể được tính toán bằng cách liệt kê các mô hình

Chúng tôi \emph{có thể} liệt kê các mô hình FOL cho từ vựng KB nhất định:

Đối với mỗi số phần tử miền \mat{$n$} từ \mat{$1$} đến \mat{$\infty$}\\
\tab Đối với mỗi vị từ \mat{$k$}-ary \mat{$P_k$} trong từ vựng\\
\tab\tab Đối với mỗi mối quan hệ-ary \mat{$k$} có thể có trên các đối tượng \mat{$n$}\\
\tab\tab\tab Đối với mỗi ký hiệu hằng số \mat{$C$} trong từ vựng\\
\tab\tab\tab\tab Cho mỗi lựa chọn tham chiếu cho \mat{$C$} từ \mat{$n$} đối tượng $\ldots$

Việc tính toán hệ quả bằng cách liệt kê các mô hình FOL không phải là điều dễ dàng!



%%%%%%%%%%%% Slide %%%%%%%%%%%%%%%%%%%%%%%%%%%%%%%%%%%%%%%%%%%%%%%%%%%
\heading{Định lượng phổ quát}

\mat{$\All{\<variables\>} \<sentence\>$}

Mọi người ở Berkeley đều thông minh:\\
\mat{$\All{x} At(x,Berkeley) \implies Smart(x)$}

\mat{$\All{x} P$}\quad đúng trong mô hình \mat{$m$} nếu \mat{$P$} đúng với \mat{$x$} là \\
\emph{mỗi} đối tượng có thể có trong mô hình

\emph{Nói đại khái là}, tương đương với liên từ \note{} của \note{insantiations} của \mat{$P$}
\mat{\[
\begin{array}{rl}
     & (At(KingJohn,Berkeley) \implies Smart(KingJohn)) \\
\land& (At(Richard,Berkeley) \implies Smart(Richard)) \\
\land& (At(Berkeley,Berkeley) \implies Smart(Berkeley)) \\
\land& \ldots
\end{array}\]}

%%%%%%%%%%%% Slide %%%%%%%%%%%%%%%%%%%%%%%%%%%%%%%%%%%%%%%%%%%%%%%%%%%
\heading{Một lỗi thường gặp cần tránh}


Thông thường, \mat{$\implies$} là từ nối chính với \mat{$\forall$}

Lỗi thường gặp: dùng \mat{$\land$} làm liên kết chính với \mat{$\forall$}:
\mat{\[\All{x} At(x,Berkeley) \land Smart(x)\]}
có nghĩa là ''Mọi người đều ở Berkeley và mọi người đều thông minh''



%%%%%%%%%%%% Slide %%%%%%%%%%%%%%%%%%%%%%%%%%%%%%%%%%%%%%%%%%%%%%%%%%%
\heading{Định lượng hiện sinh}

\mat{$\Exi{\<variables\>} \<sentence\>$}

Ai đó ở Stanford rất thông minh:\\
\mat{$\Exi{x} At(x,Stanford) \land Smart(x)$}

\mat{$\Exi{x} P$}\quad đúng trong mô hình \mat{$m$} nếu \mat{$P$} đúng với \mat{$x$} là \\
\emph{một số } đối tượng có thể có trong mô hình

\emph{Nói đại khái là}, tương đương với \note{phân cách} của \note{instantiations} của \mat{$P$}
\mat{\[
\begin{array}{rl}
     & (At(KingJohn,Stanford) \land Smart(KingJohn)) \\
\lor& (At(Richard,Stanford) \land Smart(Richard)) \\
\lor& (At(Stanford,Stanford) \land Smart(Stanford)) \\
\lor& \ldots
\end{array}\]}

%%%%%%%%%%%% Slide %%%%%%%%%%%%%%%%%%%%%%%%%%%%%%%%%%%%%%%%%%%%%%%%%%%
\heading{Một lỗi phổ biến khác cần tránh}

Thông thường, \mat{$\land$} là từ nối chính với \mat{$\exists$}

Lỗi thường gặp: dùng \mat{$\implies$} làm liên kết chính với \mat{$\exists$}:
\mat{\[\Exi{x} At(x,Stanford) \implies Smart(x)\]}
là đúng nếu có ai đó không ở Stanford!



%%%%%%%%%%%% Slide %%%%%%%%%%%%%%%%%%%%%%%%%%%%%%%%%%%%%%%%%%%%%%%%%%%
\heading{Tính chất của bộ định lượng}

\mat{$\All{x}\All{y}$} giống với \mat{$\All{y}\All{x}$} (\q{tại sao})

\mat{$\Exi{x}\Exi{y}$} giống với \mat{$\Exi{y}\Exi{x}$} (\q{tại sao})

\mat{$\Exi{x}\All{y}$} là \emph{không phải} giống với \mat{$\All{y}\Exi{x}$}

\mat{$\Exi{x}\All{y} Loves(x,y)$}\\
''Có một người trên đời yêu thương tất cả mọi người''

\mat{$\All{y}\Exi{x} Loves(x,y)$}\\
''Mọi người trên thế giới đều có ít nhất một người được yêu thương''

\note{Tính đối ngẫu định lượng}: mỗi cái có thể được biểu thị bằng cách sử dụng cái kia

\mat{$\All{x} Likes(x,IceCream)$} \qquad \mat{$\lnot \Exi{x} \lnot Likes(x,IceCream)$} 

\mat{$\Exi{x} Likes(x,Broccoli)$} \qquad \mat{$\lnot \All{x} \lnot Likes(x,Broccoli)$} 





%%%%%%%%%%%% Slide %%%%%%%%%%%%%%%%%%%%%%%%%%%%%%%%%%%%%%%%%%%%%%%%%%%
\heading{Thú vị với những câu nói}

Anh em là anh em ruột




%%%%%%%%%%%% Slide %%%%%%%%%%%%%%%%%%%%%%%%%%%%%%%%%%%%%%%%%%%%%%%%%%%
\heading{Thú vị với những câu nói}

Anh em là anh em ruột

\mat{$\All{x,y} Brother(x,y) \implies Sibling(x,y)$}.

''Anh chị em'' có tính đối xứng



%%%%%%%%%%%% Slide %%%%%%%%%%%%%%%%%%%%%%%%%%%%%%%%%%%%%%%%%%%%%%%%%%%
\heading{Thú vị với những câu nói}

Anh em là anh em ruột

\mat{$\All{x,y} Brother(x,y) \implies Sibling(x,y)$}.

''Anh chị em'' có tính đối xứng

\mat{$\All{x,y} Sibling(x,y) \lequiv Sibling(y,x)$}.

Mẹ của một người là phụ nữ của một người




%%%%%%%%%%%% Slide %%%%%%%%%%%%%%%%%%%%%%%%%%%%%%%%%%%%%%%%%%%%%%%%%%%
\heading{Thú vị với những câu nói}

Anh em là anh em ruột

\mat{$\All{x,y} Brother(x,y) \implies Sibling(x,y)$}.

''Anh chị em'' có tính đối xứng

\mat{$\All{x,y} Sibling(x,y) \lequiv Sibling(y,x)$}.

Mẹ của một người là phụ nữ của một người

\mat{$\All{x,y} Mother(x,y) \lequiv (Female(x) \land Parent(x,y))$}.

Anh em họ đầu tiên là con của anh chị em ruột của cha mẹ




%%%%%%%%%%%% Slide %%%%%%%%%%%%%%%%%%%%%%%%%%%%%%%%%%%%%%%%%%%%%%%%%%%
\heading{Thú vị với những câu nói}

Anh em là anh em ruột

\mat{$\All{x,y} Brother(x,y) \implies Sibling(x,y)$}.

''Anh chị em'' có tính đối xứng

\mat{$\All{x,y} Sibling(x,y) \lequiv Sibling(y,x)$}.

Mẹ của một người là phụ nữ của một người

\mat{$\All{x,y} Mother(x,y) \lequiv (Female(x) \land Parent(x,y))$}.

Anh em họ đầu tiên là con của anh chị em ruột của cha mẹ

\mat{$\All{x,y} FirstCousin(x,y) \lequiv 
\Exi{p,ps} Parent(p,x) \land Sibling(ps,p) \land Parent(ps,y)$}


%%%%%%%%%%%% Slide %%%%%%%%%%%%%%%%%%%%%%%%%%%%%%%%%%%%%%%%%%%%%%%%%%%
\heading{Bình đẳng}

\mat{$term_1 = term_2$} đúng theo cách giải thích nhất định\\
khi và chỉ nếu \mat{$term_1$} và \mat{$term_2$} đề cập đến cùng một đối tượng

\begin{tabular}{ll}
Ví dụ: & \mat{$1=2$} và \mat{$\All{x} {\times}(Sqrt(x),Sqrt(x)) = x$} đều thỏa mãn\\
      & \mat{$2=2$} hợp lệ
\end{tabular}

Ví dụ: định nghĩa (đầy đủ) \mat{$Sibling$} theo \mat{$Parent$}:\al
\mat{$\All{x,y} Sibling(x,y) \lequiv [\lnot(x\eq y) \land \Exi{m,f} \lnot(m\eq f) \land {}$}\nl
\mat{$Parent(m,x) \land Parent(f,x) \land Parent(m,y) \land Parent(f,y)]$}



%%%%%%%%%%%% Slide %%%%%%%%%%%%%%%%%%%%%%%%%%%%%%%%%%%%%%%%%%%%%%%%%%%
\heading{Tương tác với FOL KBs}

Giả sử một đại lý wumpus-world đang sử dụng FOL KB\\
và cảm nhận được mùi và làn gió (nhưng không lấp lánh) tại \mat{$t=5$}:

\mat{${\sc Tell}(KB,Percept([Smell,Breeze,None],5))$}\\
\mat{${\sc Ask}(KB,\Exi{a} Action(a,5))$}

Tức là, \mat{$KB$} có đòi hỏi bất kỳ hành động cụ thể nào tại \mat{$t=5$} không?

Trả lời: \mat{$Yes,\ \{a/Shoot\}$}\qquad $\leftarrow$ \defn{thay thế} (danh sách ràng buộc)

Cho một câu \mat{$S$} và một câu thay thế \mat{$\sigma$},\\
\mat{$S\sigma$} biểu thị kết quả của việc cắm \mat{$\sigma$} vào \mat{$S$}; ví dụ:\\
\mat{$S = Smarter(x,y)$}\\
\mat{$\sigma = \{x/Hillary,y/Bill\}$}\\
\mat{$S\sigma = Smarter(Hillary,Bill)$}

\mat{${\sc Ask(KB,S)}$} trả về một số/tất cả \mat{$\sigma$} sao cho \mat{$KB \models S\sigma$}



%%%%%%%%%%%% Slide %%%%%%%%%%%%%%%%%%%%%%%%%%%%%%%%%%%%%%%%%%%%%%%%%%%
\heading{Cơ sở kiến thức về thế giới wumpus}

\note{``Nhận thức''}\\
\mat{$\All{b,g,t} Percept([Smell,b,g],t) \implies Smelt(t)$}\\
\mat{$\All{s,b,t} Percept([s,b,Glitter],t) \implies AtGold(t)$}

\note{Phản xạ}: \mat{$\All{t} AtGold(t) \implies Action(Grab,t)$}

\note{Phản xạ với trạng thái bên trong}: chúng ta đã có vàng chưa?\\
\mat{$\All{t} AtGold(t) \land \lnot Holding(Gold,t) \implies Action(Grab,t)$}

\mat{$Holding(Gold,t)$} không thể quan sát được \nl
$\Rightarrow$ theo dõi sự thay đổi là điều cần thiết

%%%%%%%%%%%% Slide %%%%%%%%%%%%%%%%%%%%%%%%%%%%%%%%%%%%%%%%%%%%%%%%%%%
\heading{Suy ra thuộc tính ẩn}

Thuộc tính của vị trí:\\
\mat{$\All{x,t} At(Agent,x,t) \land Smelt(t) \implies Smelly(x)$}\\
\mat{$\All{x,t} At(Agent,x,t) \land Breeze(t) \implies Breezy(x)$}

Các quảng trường thoáng mát gần hố:

\defn{Quy tắc chẩn đoán}---suy ra nguyên nhân từ hiệu ứng\nl
      \mat{$\All{y} Breezy(y) \implies \Exi{x} Pit(x) \land Adjacent(x,y)$}

Quy tắc \defn{Nhân quả}---suy ra kết quả từ nguyên nhân\nl
      \mat{$\All{x,y} Pit(x) \land Adjacent(x,y) \implies Breezy(y)$}

Cả hai điều này đều không hoàn chỉnh --- ví dụ, quy luật nhân quả không nói
liệu những quảng trường ở xa hố có thể mát mẻ không

\defn{Định nghĩa} cho vị từ \mat{$Breezy$}:\nl
      \mat{$\All{y}  Breezy(y)  \lequiv [\Exi{x} Pit(x) \land Adjacent(x,y)]$}



%%%%%%%%%%%% Slide %%%%%%%%%%%%%%%%%%%%%%%%%%%%%%%%%%%%%%%%%%%%%%%%%%%
\heading{Theo dõi sự thay đổi}

Sự thật tồn tại trong \defn{tình huống}, thay vì vĩnh viễn\\
Ví dụ: \mat{$Holding(Gold,Now)$} thay vì chỉ \mat{$Holding(Gold)$}

\defn{Tính toán tình huống} là một cách thể hiện sự thay đổi trong FOL:\nl
   Thêm đối số tình huống vào từng vị từ không vĩnh cửu\nl
   Ví dụ: \mat{$Now$} trong \mat{$Holding(Gold,Now)$} biểu thị một tình huống

Các tình huống được kết nối bằng chức năng \mat{$Result$}\\
\mat{$Result(a,s)$} là tình huống xảy ra do thực hiện \mat{$a$} trong \mat{$s$}

\epsfxsize=0,3\textwidth
\fig{\file{figures}{situations2.ps}}



%%%%%%%%%%%% Slide %%%%%%%%%%%%%%%%%%%%%%%%%%%%%%%%%%%%%%%%%%%%%%%%%%%
\heading{Mô tả hành động I}

Tiên đề ``Hiệu ứng''---mô tả những thay đổi do hành động\\
\mat{$\All{s} AtGold(s) \implies Holding(Gold,Result(Grab,s))$}

Tiên đề ``Khung''---mô tả \emph{không thay đổi} do hành động\\
\mat{$\All{s} HaveArrow(s) \implies HaveArrow(Result(Grab,s))$}

\defn{Vấn đề về khung}: tìm một cách tinh tế để xử lý việc không thay đổi\nl
   (a) biểu diễn---tránh tiên đề khung\nl
   (b) suy luận---tránh lặp lại ``sao chép'' để theo dõi trạng thái

\defn{Vấn đề về trình độ }: mô tả chân thực về các hành động thực tế đòi hỏi phải có vô số cảnh báo---điều gì sẽ xảy ra nếu vàng trơn hoặc bị đóng đinh hoặc $\ldots$

\defn{Vấn đề phân nhánh}: hành động thực tế gây ra nhiều hậu quả thứ yếu---còn bụi trên vàng, hao mòn trên găng tay thì sao, $\ldots$


%%%%%%%%%%%% Slide %%%%%%%%%%%%%%%%%%%%%%%%%%%%%%%%%%%%%%%%%%%%%%%%%%%
\heading{Mô tả hành động II}

\defn{Các tiên đề trạng thái kế tiếp} giải quyết vấn đề về khung biểu diễn 

Mỗi tiên đề là ``about'' một \emph{vị ngữ} (bản thân nó không phải là một hành động):
\mat{\begin{eqnarray*}
\mbox{P true afterwards}&\lequiv& [\mbox{an action made P true}\\
                      &\lor   & \mbox{P true already and no action
                                  made P false}]
\end{eqnarray*}}

Để giữ vàng:\al
   \mat{$\All{a,s} Holding(Gold,Result(a,s)) \lequiv {}$}\nl
      \mat{$[(a\eq Grab \land AtGold(s))$}\nl
      \mat{${}\lor (Holding(Gold,s) \land a\neq Release)]$}


%%%%%%%%%%%% Slide %%%%%%%%%%%%%%%%%%%%%%%%%%%%%%%%%%%%%%%%%%%%%%%%%%%
\heading{Lập kế hoạch}

Điều kiện ban đầu tính bằng KB:\nl
   \mat{$At(Agent,[1,1],S_0)$}\nl
   \mat{$At(Gold,[1,2],S_0)$}

Truy vấn: \mat{${\sc Ask}(KB,\Exi{s} Holding(Gold,s))$}\nl
   tức là tôi sẽ giữ vàng trong tình huống nào?

Trả lời: \mat{$\{s/Result(Grab,Result(Forward,S_0))\}$}\nl
   tức là đi về phía trước và lấy vàng

Điều này giả định rằng đại lý quan tâm đến các kế hoạch bắt đầu từ \mat{$S_0$}
và \mat{$S_0$} là tình huống duy nhất được mô tả trong KB


%%%%%%%%%%%% Slide %%%%%%%%%%%%%%%%%%%%%%%%%%%%%%%%%%%%%%%%%%%%%%%%%%%
\heading{Lập kế hoạch: Một cách tốt hơn}

Trình bày \note{kế hoạch} dưới dạng chuỗi hành động \mat{$[a_1,a_2,\ldots,a_n]$}

\mat{$PlanResult(p,s)$} là kết quả của việc thực thi \mat{$p$} trong \mat{$s$}

Sau đó truy vấn \mat{${\sc Ask}(KB,\Exi{p} Holding(Gold,PlanResult(p,S_0)))$}\\
có giải pháp \mat{$\{p/[Forward,Grab]\}$}

Định nghĩa \mat{$PlanResult$} theo \mat{$Result$}:\al
   \mat{$\All{s} PlanResult([\,],s) = s$}\al
   \mat{$\All{a,p,s} PlanResult([a|p],s) = PlanResult(p,Result(a,s))$}

\defn{Hệ thống lập kế hoạch} là các bộ lý luận có mục đích đặc biệt được thiết kế để thực hiện việc này
loại suy luận hiệu quả hơn so với một lý do có mục đích chung

%%%%%%%%%%%% Slide %%%%%%%%%%%%%%%%%%%%%%%%%%%%%%%%%%%%%%%%%%%%%%%%%%%
\heading{Tóm tắt}

Logic bậc nhất: \al
-- các đối tượng và quan hệ là những nguyên hàm ngữ nghĩa\al
-- cú pháp: hằng, hàm, vị ngữ, đẳng thức, định lượng

Khả năng biểu cảm tăng lên: đủ để định nghĩa thế giới wumpus 

Tính toán tình huống:\al
-- quy ước mô tả hành động và thay đổi trong FOL\al
-- có thể xây dựng kế hoạch dưới dạng suy luận về một phép tính tình huống KB


\end{huge} 
\end{document}
